\section*{Capítulo \ref{ch:g12:domesticated-plants} --- La planta domesticada}

\begin{solution}{exo:g12:domesticated-plants:1}
\omterm{def:g12:domesticated-plants:domestication}{Domesticación}: la transformación de una \omterm{def:g10:biodiversity-scales:species}{especie} silvestre en otra que
depende de los seres humanos y los sirve. \omterm{def:g12:domesticated-plants:domestication}{Selección artificial}: que las
personas guarden y siembren los individuos con los caracteres que
quieren, y suban así la frecuencia de los \omterm{def:g10:universal-dna:gene}{alelos} responsables de ellos.
\end{solution}

\begin{solution}{exo:g12:domesticated-plants:2}
Espigas que no se desgranan (no hay dispersión); germinación inmediata
(no se puede esperar a que pase un mal año); pérdida de las toxinas (se
lo come todo el mundo); una sola espiga grande en una planta sin
ramificar (todos los huevos en la misma cesta); maduración simultánea
(una mala semana destruye la cosecha entera).
\end{solution}

\begin{solution}{exo:g12:domesticated-plants:3}
La descendencia uniforme y vigorosa de dos líneas puras homocigóticas.
Su propia semilla segrega en una mezcla de plantas menos vigorosas, así
que el agricultor compra semilla híbrida nueva cada año.
\end{solution}

\begin{solution}{exo:g12:domesticated-plants:4}
El \omterm{def:g10:universal-dna:gene}{gen} se inserta en el \omterm{def:g10:universal-dna:information}{ADN} que una bacteria del suelo transfiere de
forma natural a las \omterm{def:g10:cells-common-unit:cell}{células} vegetales; las \omterm{def:g10:cells-common-unit:cell}{células} infectadas que toman
el \omterm{def:g10:universal-dna:gene}{gen} se regeneran hasta dar plantas enteras.
\end{solution}

\begin{solution}{exo:g12:domesticated-plants:5}
De unas \num{1.9} a \qty{7.2}{t/ha}: casi cuatro veces.
\end{solution}

\begin{solution}{exo:g12:domesticated-plants:6}
Las plantas cuyo grano se desgranaba lo perdían antes de la cosecha o
durante ella; los mutantes de tallo resistente conservaban el suyo y
eran los que se recogían, y por tanto los que se sembraban. La cosecha
misma seleccionó el \omterm{def:g10:universal-dna:gene}{alelo} de tallo resistente, generación tras
generación.
\end{solution}

\begin{solution}{exo:g12:domesticated-plants:7}
Una \omterm{def:g10:biodiversity-scales:species}{especie}: se cruzan entre sí y dan descendencia fértil. Los
agricultores seleccionaron la exageración de un órgano distinto en cada
linaje --- hojas, yema terminal, botones florales ---, de modo que se
diferencian en la regulación del crecimiento, no en su capacidad de
cruzarse.
\end{solution}

\begin{solution}{exo:g12:domesticated-plants:8}
La F1 es toda $Aa$. La F2: $AA$ $1/4$, $Aa$ $1/2$, $aa$ $1/4$. La F1 es
uniforme porque cada planta recibe $A$ de un progenitor y $a$ del otro;
en la F2, la \omterm{def:g12:meiosis-genetic-shuffling:meiosis}{meiosis} y la fecundación reparten los \omterm{def:g10:universal-dna:gene}{alelos} al azar.
\end{solution}

\begin{solution}{exo:g12:domesticated-plants:9}
Entre millones de barrenadores, unos pocos llevaban una \omterm{def:g11:mutations:mutation}{mutación} que
los hacía tolerar la toxina; en un campo en el que todas las plantas la
fabrican, solo ellos sobrevivieron y se reprodujeron --- selección. Un
refugio mantiene vivos a muchos barrenadores sensibles; sus cruces con
los raros \omterm{prop:g11:antibiotic-resistance:mechanisms}{resistentes} dan descendientes que llevan un solo \omterm{def:g10:universal-dna:gene}{alelo} de
resistencia y que, al ser \omterm{def:g11:genes-and-disease:genetic}{recesivo}, mueren en el maíz con toxina, de
modo que el \omterm{def:g10:universal-dna:gene}{alelo} se diluye en cada generación.
\end{solution}

\begin{solution}{exo:g12:domesticated-plants:10}
Los \omterm{def:g10:universal-dna:gene}{genes} de que se trata controlan el desarrollo --- la ramificación
de la planta, el crecimiento de la cáscara del grano, el número de
filas ---; un cambio en la regulación de uno de esos \omterm{def:g10:universal-dna:gene}{genes} altera la
forma de un órgano entero. Pocos cambios reguladores, grandes efectos
visibles.
\end{solution}

\begin{solution}{exo:g12:domesticated-plants:11}
Guardan los \omterm{def:g10:universal-dna:gene}{alelos} que la selección eliminó del cultivo: resistencia a
enfermedades y plagas, tolerancia a la sequía, a la sal o al frío.
Cuando llega una enfermedad nueva o un clima nuevo, los mejoradores
devuelven esos \omterm{def:g10:universal-dna:gene}{alelos} al cultivo por cruzamiento; una vez extinguido un
pariente silvestre, sus \omterm{def:g10:universal-dna:gene}{alelos} se han perdido.
\end{solution}

\begin{solution}{exo:g12:domesticated-plants:12}
El polen de la colza fecunda a la mostaza; los híbridos llevan el \omterm{def:g10:universal-dna:gene}{gen}
de tolerancia; retrocruzados con la mostaza durante algunas
generaciones y seleccionados por el herbicida que se pulveriza en los
campos, se extiende una mostaza silvestre tolerante. Procesos:
hibridación y paso horizontal de un \omterm{def:g10:universal-dna:gene}{gen} entre \omterm{def:g12:selection-drift-speciation:population}{poblaciones}, y después
\omterm{def:g12:selection-drift-speciation:selection}{selección natural} por el herbicida.
\end{solution}

\begin{solution}{exo:g12:domesticated-plants:13}
Biológicas: la \omterm{def:g10:cell-metabolism:autotrophy}{fotosíntesis} de la planta y su capacidad de llenar el
grano están cerca de sus límites, y la \omterm{def:g10:biodiversity-scales:biodiversity}{diversidad genética} de las
variedades de élite, de la que deben salir las mejoras siguientes, es
estrecha. No biológica: el uso de fertilizantes y de plaguicidas ya no
aumenta, o se está reduciendo, y las tensiones del clima han crecido.
Un mejorador podría traer \omterm{def:g10:universal-dna:gene}{alelos} de los parientes silvestres, o editar
\omterm{def:g10:universal-dna:gene}{genes} de la \omterm{def:g10:cell-metabolism:autotrophy}{fotosíntesis} o de la tolerancia a la sequía.
\end{solution}

\begin{solution}{exo:g12:domesticated-plants:14}
Cruzamiento: el \omterm{def:g10:universal-dna:gene}{alelo} tiene que existir en una planta que pueda
cruzarse con el cultivo; se rebarajan miles de \omterm{def:g10:universal-dna:gene}{genes} que hay que
separar después por selección a lo largo de unas diez generaciones.
\omterm{prop:g10:universal-dna:universal}{Transgénesis}: el \omterm{def:g10:universal-dna:gene}{gen} puede venir de cualquier organismo; se añade un
solo \omterm{def:g10:universal-dna:gene}{gen} y el resto del genoma queda intacto; unos pocos años.
\end{solution}

\begin{solution}{exo:g12:domesticated-plants:15}
La aritmética es idéntica: los individuos que llevan ciertos \omterm{def:g10:universal-dna:gene}{alelos} se
reproducen más, y las \omterm{def:g12:selection-drift-speciation:population}{frecuencias} de esos \omterm{def:g10:universal-dna:gene}{alelos} suben. Lo que cambia
es el agente --- una elección humana en lugar del ambiente --- y el
desenlace para la planta: la \omterm{def:g12:selection-drift-speciation:selection}{selección natural} mantiene a una \omterm{def:g10:biodiversity-scales:species}{especie}
capaz de vivir por su cuenta, mientras que la \omterm{def:g12:domesticated-plants:domestication}{domesticación} selecciona
caracteres que sirven a quien selecciona y deja a la planta dependiente
de él. La \omterm{def:g12:domesticated-plants:domestication}{domesticación} es \omterm{def:g12:selection-drift-speciation:selection}{selección natural} con otro seleccionador, no
su contrario.
\end{solution}

\begin{solution}{pb:g12:domesticated-plants:1}
\textbf{1.} Teosinte $12 \times 0.08 \approx \qty{1}{g}$; maíz $600
\times 0.3 = \qty{180}{g}$: casi doscientas veces más.

\textbf{2.} Muchos granos en un zuro rígido (no hay dispersión); granos
desnudos (sin protección frente a los animales y la podredumbre); una
sola espiga grande en un único tallo (no se reparte el riesgo); granos
sujetos dentro de las brácteas (no pueden ni caer al suelo).

\textbf{3.} El \omterm{def:g10:universal-dna:gene}{alelo} desnudo es \omterm{def:g11:genes-and-disease:genetic}{recesivo}: una planta muestra granos
desnudos solo cuando lleva dos copias, lo que, para un \omterm{def:g10:universal-dna:gene}{alelo} raro,
ocurre en una minoría ínfima de plantas; los \omterm{def:g11:genes-and-disease:genetic}{portadores} de una sola
copia tienen aspecto corriente.

\textbf{4.} $q^2 = 0.0001$: una planta de cada diez mil.

\textbf{5.} Toda la semilla viene de plantas $nn$, cuyos granos
recibieron $n$ de la madre; el polen vino sobre todo del campo
corriente ($N$ con frecuencia 0,99). La semilla es casi toda $Nn$, así
que la frecuencia de $n$ salta a cerca de 0,5 --- una generación de
selección severa.

\textbf{6.} $600/12 = 50$ en 9000 generaciones: un factor de
$50^{1/9000}
\approx 1.0004$ por generación, una ganancia del 0,04\%.
Pasos ínfimos, compuestos a lo largo de miles de generaciones, se
convierten en un cambio grande --- la aritmética de la selección.

\textbf{7.} Con muchos \omterm{def:g10:universal-dna:gene}{genes} que aportan cada uno un poco, la selección
sube la frecuencia de los \omterm{def:g10:universal-dna:gene}{alelos} favorables en todos ellos y sigue
encontrando combinaciones nuevas durante mucho tiempo; un carácter de
un solo \omterm{def:g10:universal-dna:gene}{gen} cambia de golpe en cuanto el \omterm{def:g10:universal-dna:gene}{alelo} \omterm{def:g11:genes-and-disease:genetic}{recesivo} se vuelve
homocigótico.

\textbf{8.} Una \omterm{def:g10:biodiversity-scales:species}{especie}: los híbridos son fértiles. La \omterm{def:g12:domesticated-plants:domestication}{domesticación}
cambió las \omterm{def:g12:selection-drift-speciation:population}{frecuencias} de los \omterm{def:g10:universal-dna:gene}{alelos} de unos pocos \omterm{def:g10:universal-dna:gene}{genes} y la forma de
la planta; no creó ninguna barrera reproductiva.

\textbf{9.} Una enfermedad o una plaga que venza a una planta las vence
a todas; la respuesta es la diversidad guardada en las variedades
antiguas y en los parientes silvestres, en bancos de semillas y en el
campo.

\textbf{10.} El agricultor es la dispersión de la planta, su ambiente
selectivo y su única vía hacia la generación siguiente: el maíz está
sometido a \omterm{def:g12:domesticated-plants:domestication}{selección artificial} en cada generación y no puede salir de
ella.

\textbf{11.} $10/5.5 - 1 \approx 82\%$ por encima del mejor de los
padres; $1 - 8.2/10 = 18\%$ por debajo de la F1.

\textbf{12.} Cada línea es homocigótica para algunos \omterm{def:g10:universal-dna:gene}{alelos} \omterm{def:g11:genes-and-disease:genetic}{recesivos}
débiles; en la F1, cada uno de esos \omterm{def:g10:universal-dna:gene}{alelos} queda enmascarado por el
\omterm{def:g10:universal-dna:gene}{alelo} funcional de la otra línea. En la F2, una cuarta parte de las
plantas es homocigótica para el \omterm{def:g10:universal-dna:gene}{alelo} débil en cada uno de esos \omterm{def:g10:universal-dna:gene}{genes},
y la media baja.

\textbf{13.} Comprando: $10 - 0.4 = \qty{9.6}{t/ha}$; guardando semilla:
\qty{8.2}{t/ha}. Comprar gana por \qty{1.4}{t/ha} cada año, y por eso
una empresa que posee las líneas puras puede vender la semilla híbrida
todas las temporadas.

\textbf{14.} La empresa se queda con las líneas; con ellas, un
agricultor podría hacer el híbrido indefinidamente y la empresa
vendería una sola vez. Las líneas son el capital de la empresa; solo se
vende su producto.

\textbf{15.} Todos los campos de una variedad híbrida son genéticamente
idénticos, y las regiones cultivan unos pocos híbridos: la diversidad
se estrecha. El agricultor ya no produce semilla y depende del
proveedor cada año.

\textbf{16.} $q^2 = 10^{-6}$: sobrevive un barrenador de cada millón;
entre los supervivientes, todos $rr$, la frecuencia de $r$ es 1.

\textbf{17.} Los supervivientes se cruzan entre sí; su descendencia es
toda $rr$; en unas pocas generaciones, los barrenadores del campo son
\omterm{prop:g11:antibiotic-resistance:mechanisms}{resistentes} --- como las bacterias de un paciente que no completa el
tratamiento.

\textbf{18.} Un superviviente \omterm{prop:g11:antibiotic-resistance:mechanisms}{resistente} ($rr$) que se cruza con un
barrenador del refugio ($RR$, la inmensa mayoría) da descendencia $Rr$,
que muere en el maíz con toxina. Mientras los barrenadores sensibles
sean muchísimo más numerosos que los \omterm{prop:g11:antibiotic-resistance:mechanisms}{resistentes}, casi todos los hijos
de un barrenador \omterm{prop:g11:antibiotic-resistance:mechanisms}{resistente} son heterocigóticos y se eliminan, y $r$
sigue siendo raro.

\textbf{19.} En qué se parecen: los dos buscan impedir la
multiplicación de los raros individuos \omterm{prop:g11:antibiotic-resistance:mechanisms}{resistentes}. En qué se
diferencian: el tratamiento completo los elimina matando a todos los
individuos antes de que la resistencia pueda seleccionarse; el refugio
acepta supervivientes y los diluye con parejas sensibles.

\textbf{20.} Casi doscientas veces más grano por espiga; unos cinco
\omterm{def:g10:universal-dna:gene}{genes} principales; una planta que depende del agricultor para
dispersarse y sembrarse.
\end{solution}
